\section{Introduction}\label{sec:introduction}

An Egyptian-fraction expansion of a positive rational number is a finite
sum of distinct reciprocals of positive integers. For integers
$1\le a<b$, let $N(a,b)$ be the least $k$ for which
\[
 \frac ab=\frac1{n_1}+\cdots+\frac1{n_k},
 \qquad 2\le n_1<\cdots<n_k,\qquad n_i\in\Z,
\]
and put $N(b)=\max_{1\le a<b}N(a,b)$. The fractions $a/b$ need not be
in lowest terms, and the denominators $n_i$ have no prescribed upper bound.
The greedy algorithm shows that such expansions exist. We determine
the order of their maximum minimum length. Individual numerators can
have much shorter expansions: for example, $N(1,b)=1$ for every $b\ge2$.

\begin{theorem}\label{thm:main}
There are absolute constants $c_1,c_2>0$ and $b_0$ such that, for every
integer $b\ge b_0$,
\[
 c_1\log\log b\le N(b)\le c_2\log\log b.
\]
In particular, the upper bound holds for every integer numerator
$1\le a<b$.
\end{theorem}

Nakayama's study of $N(a,b)$ emphasized arithmetic criteria for
expansions with few terms \cite[Section~I]{Nakayama1940}.
The uniform problem asks how these minimum lengths behave as the
numerator varies over a fixed denominator. In 1950, Erd\H{o}s
reported de Bruijn's bound $N(b)\ll\log b/\log\log\log b$ and
improved it to $N(b)\ll\log b/\log\log b$
\cite[p.~195, Theorem~1]{Erdos1950}.
He proposed the double-logarithmic upper bound and proved the matching
lower order, already for the numerator $b-1$
\cite[p.~195, Theorem~2]{Erdos1950}.
The question appears again in Erd\H{o}s and Graham
\cite[pp.~37--38]{ErdosGraham1980} and is recorded as Erd\H{o}s
Problem~304 \cite{Bloom304}.
Vose subsequently obtained $N(b)\ll\sqrt{\log b}$ \cite{Vose1985};
an explicit modern statement of his uniform bound is given in
\cite[Lemma~2.2]{vanDoornTang2026}.
Theorem~\ref{thm:main} proves Erd\H{o}s's conjecture affirmatively.
The lower bound is classical; the issue is the uniform upper bound.

A related line of work controls the denominators as well as the
length. Tenenbaum and Yokota proved that, for each fixed
$\varepsilon>0$ and every sufficiently large $b$, every $a/b\in(0,1)$
has a distinct expansion of length at most
$(1+\varepsilon)\log b/\log\log b$, with all denominators at most
$4b(\log b)^2\log\log b$ \cite[Theorem, p.~151]{TenenbaumYokota1990}.
Their length estimate serves this simultaneous constraint.
Here we optimize only the order of the length and leave the
denominator sizes unrestricted.

The uniform theorem also determines the double-logarithmic order of
the number of representations of one. For each positive integer $k$, put
\[
 F(k)=\#\left\{(n_1,\ldots,n_k):
  1\le n_1<\cdots<n_k,\quad n_i\in\Z,\quad
  \frac1{n_1}+\cdots+\frac1{n_k}=1\right\}.
\]
Here too there is no upper bound on the denominators.

\begin{corollary}\label{cor:counting}
There are absolute constants $c,C>0$ and $k_0$ such that, for every
integer $k\ge k_0$,
\[
 ck\le\log\log F(k)\le Ck.
\]
\end{corollary}

Erd\H{o}s and Graham asked for good estimates for this counting function
\cite[p.~32]{ErdosGraham1980}. Konyagin stated a lower bound with
double logarithm of order $k/\log k$ \cite[Theorem~1]{Konyagin2014};
Elsholtz obtained this lower order even when every denominator is
congruent to $1$ or $-1$ modulo a fixed squarefree integer $P>1$:
the bound holds for every sufficiently large $k$, with $k$ required
to be odd when $P$ is even, and its constant may depend on $P$
\cite[Theorem~1.1]{Elsholtz2016}. Elsholtz and Planitzer proved an upper
bound with double logarithm $O(k)$, allowing repetitions as well
\cite[Corollary~3(2)]{ElsholtzPlanitzer2021}.
Corollary~\ref{cor:counting} determines the double-logarithmic order for
unrestricted distinct denominators.


We can also ask which exact denominators occur among these expansions.
For each integer $k\ge1$, let $D_k$ be the set of integers $m\ge2$
for which there are positive integers $n_1<\cdots<n_k$ satisfying
\[
 1=\sum_{i=1}^{k}\frac1{n_i},\qquad
 m\in\{n_1,\ldots,n_k\}.
\]
Set
\[
 v(k)=\min\bigl(\{2,3,\ldots\}\setminus D_k\bigr).
\]
The elementary denominator bound proved below makes $D_k$ finite, so
this minimum exists. Each $m\in D_k$ may use a different expansion;
membership requires the exact denominator $m$ and exactly $k$ distinct
terms.

\begin{corollary}\label{cor:prescribed}
For every sufficiently large integer $k$,
\[
 \exp\bigl(\exp(k/600)\bigr)\le v(k)\le1+k^{2^{k-1}}.
\]
More precisely,
\[
 \frac{\log2}{257}
 \le\liminf_{k\to\infty}\frac{\log\log v(k)}k
 \le\limsup_{k\to\infty}\frac{\log\log v(k)}k
 \le\log2.
\]
Equivalently for the lower bound, every fixed $c<\log2/257$ satisfies
$v(k)\ge\exp(\exp(ck))$ for every sufficiently large integer $k$.
In particular $\log\log v(k)=\Theta(k)$.
\end{corollary}

The prescribed-denominator question originates in Erd\H{o}s and Graham
\cite[p.~35]{ErdosGraham1980}; see also Problem~293 \cite{Bloom293}.
Van Doorn and Tang prove the lower bound $v(k)\ge\exp(c k^2)$ for
an absolute $c>0$ \cite[Theorem~1.1]{vanDoornTang2026}. They also show
that an exact prescribed denominator can be
retained while increasing the length of a distinct expansion
\cite[Lemma~2.1]{vanDoornTang2026} and anticipate the connection
between uniform short expansions and a double-exponential lower bound
for $v(k)$ \cite[Section~3]{vanDoornTang2026}. We give the required
marker-avoidance argument below, thereby deriving the qualitative order
from Theorem~\ref{thm:main}. A direct construction additionally gives
the numerical lower slope in Corollary~\ref{cor:prescribed}, without
quantifying the constant $c_2$ in the uniform theorem. The elementary
upper bound is convenient rather than best known: sharper public upper
bounds follow from the counting estimates of Elsholtz and Planitzer
\cite[Corollary~3(2)]{ElsholtzPlanitzer2021}.

The two consequences impose different requirements on a construction.
To obtain many expansions, it suffices to retain a denominator divisible
by an integer with many prime factors. To include a prescribed integer
$m$ in $D_k$, the exact term $1/m$ must survive. The latter requirement
governs both the initial construction and the operation that increases
its length.

\subsection*{The proof in outline}

Put $S=\log b$. After $O(\log S)$ greedy steps, either the expansion
is complete or its remainder has the form $A/C$, where $A\le b$ and
$C$ lies in a prescribed exponential range in $S$.
We then construct an auxiliary integer $M$ with $e^S<M=\exp(O(S))$
and a large interval of numerators in which almost every fraction $u/(MC)$ has
an expansion of length $O(\log S)$. This dense family suffices for
the original numerator: a suitable integer multiple $gAM$ lies well
inside the interval, and it is the sum of two numerators from that dense family.
Adding their expansions and dividing by $g$ represents $A/C$.

The dense family is obtained by descending through $O(\log S)$
numerator ranges. At a range whose upper endpoint exceeds $e^S$, we
represent fractions $u/(MC)$; at lower ranges we represent $u/M$,
which is still less than one. Write $Q=C$ or $Q=1$ for the denominator
factor in the current range. A divisor $t$ of $M$ gives the identity
\[
 \frac{u}{MQ}
 =\frac1{(M/t)z}+\frac1z\frac{h}{MQ},
 \qquad z=\left\lceil\frac{Qt}{u}\right\rceil,
 \qquad h=uz-Qt.
\]
Thus one unit fraction reduces the problem to the residue of $-Qt$
modulo $u$. When the denominator factor changes from $C$ to $1$, an
expansion over $M$ can be transferred back over $MC$ by multiplying
every denominator by $C$. We choose $M$ so that, for most numerators
at each range, many of its divisors give small residues. Deterministic products of
distinct primes provide the large-range estimate; independent random
products provide the small-range estimate. At the bottom, the same
construction gives a reduction for every remaining numerator, followed
by a binary expansion.

The difficulty is to prevent the exceptional numerators from accumulating
as these reductions are joined. For fixed $t$ and $h$, a predecessor
$u$ must divide $Qt+h$. A uniform moment estimate for the number of
small divisors in a shifted interval controls the total number of
predecessors of a small exceptional set. H\"older's inequality then gives
a recurrence for the exceptional proportions that remains small throughout the descent.
This combination of residue distribution, uniform divisor moments,
and the final passage from a dense family to every numerator is the
main mechanism of the proof.

The use of a common auxiliary denominator connects the proof to the
classical divisor method. Erd\H{o}s's construction writes suitable
integers as sums of distinct divisors of a factorial and turns those
sums into unit fractions \cite[Section~1, pp.~199--203]{Erdos1950}.
Tenenbaum and Yokota obtain a shorter divisor decomposition by choosing
successive divisors close to the remaining numerator
\cite[Lemma~4 and Section~3]{TenenbaumYokota1990}.
The present descent instead selects divisors through the residues
they produce. Modular control by finite Fourier methods also occurs
in Croot's work \cite{Croot1999} and Martin's development of it
\cite[Section~4, Lemmas~11--13]{Martin2000}, where subset sums of
modular inverses remove large prime-power denominator factors.
Here prime-product distributions supply many possible reductions,
and a uniform divisor moment keeps their exceptional sets under
control across all levels. The moment proof uses the prime-factor
splitting method of Erd\H{o}s \cite[Section~3]{Erdos1952Divisors};
the truncation and shift uniformity needed here are established below.

Section~\ref{sec:elementary} isolates the dense-family statement and
deduces Theorem~\ref{thm:main} from it. Section~\ref{sec:divisors}
proves the divisor estimate, Section~\ref{sec:residues} constructs
the residue distributions, and Section~\ref{sec:descent} completes
the dense-family argument. The construction constants are fixed first,
then a moment order large enough for the descent, and finally a
sufficiently large lower bound on $S$.
Section~\ref{sec:counting} then deduces Corollary~\ref{cor:counting}.
Theorem~\ref{thm:main} supplies a short expansion containing a denominator
with many prime factors. Splitting that denominator over its divisors
gives many distinct expansions, and an injective padding operation gives
every sufficiently large prescribed length. Divisor-rich denominators
and divisor-indexed splitting appear in Konyagin and Elsholtz
\cite{Konyagin2014,Elsholtz2016}; Konyagin explicitly records the padding
injection \cite[p.~312]{Konyagin2014}.


For the prescribed-denominator consequence, Section~\ref{sec:prescribed}
reserves $1/m$ and uses a greedy prefix that skips denominator $m$.
The remaining sum is smaller than every reserved reciprocal, so any
positive expansion of that remainder avoids all reserved denominators.
Applying Theorem~\ref{thm:main} gives a short marked expansion, and a
separate padding lemma retains the exact marker at every later length.
For an explicit length constant, Section~\ref{sec:prescribed-quantitative}
instead constructs the tail from a common supply of rational numbers
whose numerators divide one auxiliary integer. Prime-product residue
collisions and the quantitative three-prime theorem provide this supply.
A divisor-density property of the unreduced greedy denominator then
groups the remaining numerator into few pieces. The resulting length
coefficient $257/\log2$ gives the lower slope after padding and absorbing
the finitely many small markers.

Throughout, $\N=\{1,2,\ldots\}$, $\log$ is the natural logarithm, $\log_2 x=(\log x)/(\log 2)$,
$\e(x)=\exp(2\pi i x)$, and sums over real intervals have integer
indices. The notations $O(\cdot)$ and $\ll$ have absolute implied
constants unless a dependence is indicated. Intermediate unit-fraction
lists may contain repetitions. Lemma~\ref{lem:distinct} removes them
when the total is below one; Section~\ref{sec:counting} supplies the
argument at total one.
